\chapter{Git and GitHub: Version Control for AI-Assisted Projects}
\label{app:git}

\begin{remark}[Learning Objectives]
After working through this appendix, the reader should be able to:
\begin{enumerate}
  \item Explain why version control is especially important when working with
        AI coding assistants, and describe the specific failure modes it mitigates.
  \item Initialise a Git repository, stage and commit changes, and interpret
        the output of \texttt{git status}, \texttt{git log}, and \texttt{git diff}.
  \item Use branching and merging to isolate AI-generated work from stable code,
        and recover from unwanted AI changes using \texttt{git checkout} and
        \texttt{git revert}.
  \item Push a repository to GitHub, open a pull request, and link AI commit
        history to project milestones.
  \item Apply a disciplined commit-per-step workflow that makes AI contributions
        reviewable, attributable, and reversible.
\end{enumerate}
\end{remark}

% ---------------------------------------------------------------
\section{Why Version Control Matters for AI-Assisted Work}
\label{app:git:why}

\begin{context}
AI coding assistants---Claude Code, Cline, GitHub Copilot, and their
successors---share a property that distinguishes them from human collaborators:
they can generate large volumes of code changes in seconds, across many files
simultaneously, with plausible-looking output that may nonetheless contain
subtle errors.  A traditional code review catches one human's typo or logic
error across a handful of files.  An AI session might touch fifty files in
three minutes.  Without version control, reviewing, attributing, and reversing
those changes is nearly impossible.

Version control systems record the \emph{history} of every change to every
file in a project.  Git, the dominant distributed version control system,
does this by storing a directed acyclic graph of \emph{commits}---snapshots
of the project at a point in time---linked by parent pointers.  Each commit
carries a message, an author, a timestamp, and a cryptographic hash of the
entire project state.  The result is an audit log: one can determine exactly
what changed, when, and in which session---including AI-assisted sessions.

This audit log serves four functions when working with AI:

\begin{enumerate}
  \item \textbf{Reversibility.}  If an AI agent introduces a regression---a
        previously passing test now fails, or a file is unintentionally
        overwritten---\texttt{git revert} or \texttt{git reset} can restore
        any prior state in seconds.  Without Git, recovery requires manual
        reconstruction.

  \item \textbf{Attribution.}  Knowing that a specific function was written
        by an AI in session~7 of a project, rather than by a human collaborator,
        carries interpretive weight during code review.  AI-generated code
        deserves additional scrutiny for hallucinated API calls, incorrect
        assumptions about library versions, and subtle logic errors that read
        fluently but execute incorrectly.

  \item \textbf{Incremental review.}  A commit containing ten lines of change
        is reviewable in minutes.  A commit containing two thousand lines across
        thirty files is not.  Disciplined small commits---one per completed
        sub-task---keep AI contributions within human review bandwidth.

  \item \textbf{Collaboration.}  Multiple researchers working in parallel on
        different parts of a project, some using AI assistants and some not,
        can merge their work without conflict using Git's branch-and-merge
        model.
\end{enumerate}

\begin{remark}[Git is not a backup system]
\label{rem:git-not-backup}
Git tracks changes to text files efficiently but is not designed for large
binary files (datasets, compiled models, PDF outputs).  For research projects,
keep \texttt{*.csv}, \texttt{*.parquet}, \texttt{*.pkl}, and \texttt{*.h5}
files out of the repository by listing them in \texttt{.gitignore}, and store
them separately in a data store (institutional storage, S3, DVC).  This keeps
repositories fast and readable.
\end{remark}
\end{context}

% ---------------------------------------------------------------
\section{Installation and First Repository}
\label{app:git:install}

\subsection{Installing Git}
\label{app:git:install-git}

Git is pre-installed on macOS (via Xcode Command Line Tools) and most Linux
distributions.  On Windows, download Git for Windows from
\url{https://git-scm.com} or install via \texttt{winget}:

\begin{minted}{bash}
winget install --id Git.Git -e --source winget
\end{minted}

Verify the installation:

\begin{minted}{bash}
git --version
\end{minted}

Configure your identity---every commit is tagged with these values:

\begin{minted}{bash}
git config --global user.name  "Jane Smith"
git config --global user.email "jane.smith@finance-lab.edu"
\end{minted}

\subsection{Initialising a Repository}
\label{app:git:init}

To place an existing project directory under version control:

\begin{minted}{bash}
cd my-finance-project
git init
\end{minted}

This creates a hidden \texttt{.git/} directory that stores the entire history.
The working directory is unchanged.

For a new project, it is more common to create the repository on GitHub first
and then clone it locally (Section~\ref{app:git:github}).

% ---------------------------------------------------------------
\section{The Core Workflow: Stage, Commit, Push}
\label{app:git:workflow}

\begin{context}
The fundamental Git loop is:

\begin{enumerate}
  \item \textbf{Work.}  Edit files in the working directory---or let an AI
        assistant edit them.
  \item \textbf{Stage.}  Select the changes you want to record with
        \texttt{git add}.
  \item \textbf{Commit.}  Snapshot the staged changes with a descriptive
        message using \texttt{git commit}.
  \item \textbf{Push.}  Upload commits to a remote (typically GitHub) with
        \texttt{git push}.
\end{enumerate}

This cycle should be tight when working with AI.  The recommended discipline
is: \emph{commit after every completed sub-task}, not after every session.
If an AI agent refactors a function, add and commit before asking it to do
the next task.  This makes each commit small, reviewable, and reversible.
\end{context}

\subsection{Checking Status and Inspecting Changes}
\label{app:git:status}

\begin{minted}{bash}
git status           # which files have changed since last commit?
git diff             # show unstaged line-level changes
git diff --staged    # show staged line-level changes
\end{minted}

The output of \texttt{git diff} is a \emph{unified diff}: lines prefixed with
\texttt{-} were removed, lines prefixed with \texttt{+} were added.  This is
the primary interface for reviewing AI-generated changes before committing them.

\subsection{Staging and Committing}
\label{app:git:commit}

Stage specific files (preferred) or all changed files:

\begin{minted}{bash}
git add src/sentiment.py tests/test_sentiment.py   # specific files
git add .                                           # all changes (use with care)
\end{minted}

Commit with a message that explains \emph{why} the change was made, not just
what:

\begin{minted}{bash}
git commit -m "feat: replace regex tokeniser with spaCy for sentence splitting"
\end{minted}

A widely adopted convention for commit message prefixes is:

\begin{center}
\begin{tabular}{ll}
\toprule
Prefix & Meaning \\
\midrule
\texttt{feat:}   & New feature or capability \\
\texttt{fix:}    & Bug fix \\
\texttt{refine:} & Improvement without new functionality \\
\texttt{chore:}  & Maintenance (dependencies, config) \\
\texttt{docs:}   & Documentation only \\
\texttt{test:}   & Tests only \\
\bottomrule
\end{tabular}
\end{center}

When using AI assistants, add a \texttt{Co-Authored-By} trailer to the commit
message to attribute the AI contribution:

\begin{minted}{bash}
git commit -m "feat: add FinBERT sentiment pipeline

Co-Authored-By: Claude Sonnet 4.6 <noreply@anthropic.com>"
\end{minted}

\subsection{Viewing History}
\label{app:git:log}

\begin{minted}{bash}
git log --oneline          # compact one-line summary per commit
git log --stat             # show files changed per commit
git log -p                 # show full diff for each commit
git log --author="Claude"  # filter by author (or Co-Authored-By)
\end{minted}

Table~\ref{tab:git-log-example} shows a representative \texttt{git log --oneline}
output from a finance research project that used an AI assistant.

\begin{table}[ht]
\centering
\caption{Example \texttt{git log} output from an AI-assisted finance project.
Each short hash identifies a unique commit; messages distinguish human-authored
from AI-assisted work.}
\label{tab:git-log-example}
\begin{tabular}{ll}
\toprule
Hash    & Message \\
\midrule
\texttt{a3f7c2e} & feat: add FactSet API integration (human) \\
\texttt{9b1d04a} & feat: extend sentiment model to earnings calls (AI-assisted) \\
\texttt{3e8c17f} & fix: correct EDGAR filing date parsing for 10-Qs (AI-assisted) \\
\texttt{c2a5b89} & chore: update requirements.txt \\
\texttt{7f3d12e} & feat: initial FinBERT pipeline scaffold (AI-assisted) \\
\texttt{b0f9e3c} & chore: git init, add .gitignore \\
\bottomrule
\end{tabular}
\end{table}

% ---------------------------------------------------------------
\section{Undoing Changes}
\label{app:git:undo}

\begin{context}
Undoing changes is one of the most important capabilities for AI-assisted
development.  AI agents occasionally produce code that passes a superficial
review but fails in production, or silently introduces incorrect logic.  The
ability to return to a known good state---without manual reconstruction---is
what makes it safe to let an AI agent work autonomously for an extended period.
\end{context}

\subsection{Before Staging: Discarding Working-Directory Changes}

\begin{minted}{bash}
git restore path/to/file.py          # discard one file's changes
git restore .                        # discard all unstaged changes
\end{minted}

These commands are \emph{destructive}: unsaved working-directory changes are
gone permanently.  Use them when an AI has produced a bad edit and you want
to start over.

\subsection{After Committing: Reverting and Resetting}

\begin{minted}{bash}
# Safe: create a new commit that undoes the changes of commit <hash>
git revert <hash>

# More powerful: move HEAD back to <hash>, discarding later commits
# --soft keeps changes staged; --mixed keeps them unstaged; --hard discards them
git reset --hard <hash>
\end{minted}

Prefer \texttt{git revert} for work that has been pushed to a shared remote,
because it adds to history without rewriting it.  Use \texttt{git reset} only
on local-only commits.

\begin{remark}[Checkpoint discipline with Claude Code]
\label{rem:checkpoint-claude}
Claude Code takes a Git snapshot before and after every tool call that writes
a file.  These checkpoints are stored in a shadow repository inside the
session directory and can be browsed with \texttt{/diff} and restored with
\texttt{/revert} within a session.  However, they are \emph{not} visible to
ordinary \texttt{git log} commands on the project repository.  The recommended
practice is to commit to the project repository at meaningful boundaries---not
to rely on Claude Code's internal checkpoints as the sole audit trail.
\end{remark}

% ---------------------------------------------------------------
\section{Branches}
\label{app:git:branches}

A \emph{branch} is a named pointer to a commit.  Creating a branch is
instantaneous and costs nothing; it is the standard mechanism for isolating
experimental or AI-generated work from stable code.

\begin{minted}{bash}
git branch ai/edgar-sentiment   # create a branch
git switch ai/edgar-sentiment   # move to it (also: git checkout)
# ... let the AI work here ...
git switch main                 # return to main
git merge ai/edgar-sentiment    # merge if the work is good
git branch -d ai/edgar-sentiment   # delete the branch
\end{minted}

The naming convention \texttt{ai/<task>} clearly marks branches that contain
primarily AI-generated work and flags them for extra scrutiny during review.

If a merge produces \emph{conflicts}---both the branch and main modified the
same lines---Git marks the file with conflict markers:

\begin{minted}{text}
<<<<<<< HEAD
return df.groupby("ticker").mean()          # your version
=======
return df.groupby("ticker").mean(numeric_only=True)  # AI version
>>>>>>> ai/edgar-sentiment
\end{minted}

Edit the file to the correct result, remove the markers, stage the file, and
complete the merge with \texttt{git commit}.

% ---------------------------------------------------------------
\section{GitHub: Remote Repositories and Collaboration}
\label{app:git:github}

GitHub (\url{https://github.com}) hosts Git repositories in the cloud and
adds collaboration features: pull requests, code review, issue tracking, and
CI/CD integration.

\subsection{Creating a Repository on GitHub}

\begin{enumerate}
  \item Log into GitHub and click \textbf{New repository}.
  \item Choose a name, select \textbf{Private} for unpublished research, and
        optionally add a \texttt{README.md} and \texttt{.gitignore}.
  \item Copy the HTTPS or SSH remote URL.
\end{enumerate}

Link the local repository to the remote and push:

\begin{minted}{bash}
git remote add origin https://github.com/username/finance-project.git
git push -u origin main
\end{minted}

Subsequently, \texttt{git push} and \texttt{git pull} require no further
arguments (the tracking branch is set by \texttt{-u}).

\subsection{Pull Requests}
\label{app:git:pr}

A \emph{pull request} (PR) is a GitHub interface for proposing that a branch
be merged into another.  For AI-assisted work, the PR is the natural review
checkpoint:

\begin{enumerate}
  \item Push the AI-work branch to GitHub: \texttt{git push -u origin ai/edgar-sentiment}.
  \item Open a PR on GitHub: \textbf{Compare \& pull request} → write a
        description of what the AI produced and what you verified.
  \item Review the diff: every file touched by the AI is visible side-by-side
        with the previous version.  Add line-level comments where behaviour is
        unclear.
  \item Merge when satisfied; delete the branch.
\end{enumerate}

The PR thread becomes a permanent record of what the AI produced, what questions
were raised during review, and what changes were requested before merge.

\subsection{The .gitignore File}
\label{app:git:gitignore}

The \texttt{.gitignore} file tells Git to ignore matching paths.  A sensible
starting template for a Python finance project:

\begin{minted}{text}
# Python artefacts
__pycache__/
*.pyc
*.pyo
.venv/
*.egg-info/

# Jupyter
.ipynb_checkpoints/

# Data (store externally)
data/raw/
data/processed/
*.csv
*.parquet
*.h5

# API keys and secrets
.env
secrets.toml
*.key

# IDE
.vscode/settings.json
.idea/

# OS
.DS_Store
Thumbs.db
\end{minted}

Never commit files containing API keys, passwords, or proprietary data.
Use environment variables or a secrets manager instead.

% ---------------------------------------------------------------
\section{A Disciplined Workflow for AI-Assisted Finance Research}
\label{app:git:discipline}

\begin{context}
The following workflow synthesises the preceding sections into a repeatable
practice for using AI coding assistants on research projects.
\end{context}

\begin{enumerate}
  \item \textbf{Start on a branch.}  Before starting an AI session, create a
        branch: \texttt{git switch -c ai/\textit{task-name}}.  The branch name
        documents the purpose of the session.

  \item \textbf{State the task clearly.}  Write the task description in the
        AI prompt and also in a commit message or issue.  A clear record of
        what the AI was asked to do is as important as a record of what it did.

  \item \textbf{Commit after each sub-task.}  When the AI completes a
        self-contained piece of work---a function, a test, a data pipeline
        step---stop, review the diff (\texttt{git diff}), and commit.
        Do not let uncommitted AI changes accumulate across long sessions.

  \item \textbf{Review before merging.}  Open a PR even when working alone.
        The GitHub diff view is the most effective format for reviewing AI-generated
        code because it enforces line-by-line attention.

  \item \textbf{Tag stable milestones.}  Use \texttt{git tag v1.0} to mark
        commits that correspond to submitted paper drafts, reproducibility
        checkpoints, or data collection cut-offs.  Tags survive branch deletion
        and provide unambiguous version references for citations and replication.

  \item \textbf{Document the AI tools used.}  Add a short \texttt{AI.md} file
        to the repository noting which AI tools were used, for which tasks, and
        during which date ranges.  As AI tool capabilities evolve rapidly,
        future readers will benefit from knowing which version of which tool
        produced which parts of the codebase.
\end{enumerate}

\begin{exercise}[Initialising a versioned finance project]
\label{ex:git-b}
\textbf{[B]} Create a new directory called \texttt{edgar-sentiment-lab}.
Initialise a Git repository, add a Python script that downloads the most recent
10-K filing for one company from EDGAR, commit it with a meaningful message,
and push it to a new private GitHub repository.  Then create a branch
\texttt{ai/add-sentiment}, use an AI assistant to add FinBERT-based sentiment
scoring, commit the AI-generated code, and open a pull request.  Write a
two-sentence PR description explaining what the AI produced and what you
manually verified before merging.
\end{exercise}

\begin{exercise}[Recovering from a bad AI edit]
\label{ex:git-i}
\textbf{[I]} Suppose an AI assistant has modified \texttt{data\_loader.py}
and \texttt{models.py} in a single session, but you notice that the
\texttt{models.py} changes broke an existing unit test while the
\texttt{data\_loader.py} changes are correct.  Using only Git commands,
describe step by step how you would: (a)~commit the \texttt{data\_loader.py}
changes; (b)~discard the \texttt{models.py} changes and restore the file to
its pre-session state; and (c)~re-run the AI session with an explicit
instruction to fix only the test-breaking issue.
\end{exercise}

\begin{exercise}[Audit trail for reproducible research]
\label{ex:git-a}
\textbf{[A]} A finance paper uses AI-generated code to construct a
sentiment-weighted momentum signal from 20 years of earnings call transcripts.
Design a Git workflow that: (a)~distinguishes AI-assisted from human-authored
commits; (b)~links each commit to the data snapshot it was run against
(using Git tags or a \texttt{data\_manifest.json} committed alongside the
code); (c)~enables a reader to reproduce the signal by checking out a specific
tag and following a documented sequence of commands; and (d)~ensures that
no proprietary data or API keys are present in the repository history.
Describe the \texttt{.gitignore} entries, branch naming conventions, and
tagging scheme you would use.
\end{exercise}
